\documentclass[12pt]{article}
\usepackage[utf8]{inputenc}
\usepackage[T1]{fontenc}
\usepackage{hyperref}
\title{Katawixi Language Resources for K-12 Curriculum Development}
\author{Compiled by Claude Sonnet 5.5 (claude-sonnet-5-5) for Andr\'as Kornai}
\date{}
\begin{document}
\maketitle

\section{Introduction}
Katawixi (Catawishi; ISO 639-3 xat; Glottocode kata1270) is one of two Katukinan languages, spoken in Amazonas, Brazil. It is known only from two partially divergent word lists, with traditional numerals recorded only for 1 to 3. One source estimates 10 native speakers in 2006. A nearby uncontacted group may include surviving Katawixi speakers, but this is unconfirmed, and it is not known whether they speak the language of the word lists. No school teaching or curriculum exists. Contact with these people is not established, and any outside effort could harm them.

\section{Primary Educational Resources}
None found.

\section{Linguistic and Reference Materials}
\begin{itemize}
\item Wikipedia, ``Katawixi language''. \url{https://en.wikipedia.org/wiki/Katawixi_language}
\item Wikipedia, ``Katukinan languages'' (family overview). \url{https://en.wikipedia.org/wiki/Katukinan_languages}
\item Max Planck EVA numerals database entry for Katawixi. \url{https://lingweb.eva.mpg.de/channumerals/Katawixi.htm}
\item Documentation of Kanamari (ELDP archive), the other Katukinan language, listed for comparison only. \url{https://eldp.access.preservica.com/dk0158/}
\end{itemize}

\section{Religious and Community Resources}
None found. The neighbouring Hi-Merim\~a people live along the Pinhu\~a River, but no source links them to Katawixi.
\begin{itemize}
\item Wikipedia, ``Himarim\~a''. \url{https://en.wikipedia.org/wiki/Himarim%C3%A3}
\end{itemize}

\section{Digital and Audio Resources}
No audio or digital resources for Katawixi were found.

\section{Curriculum Development Context}
There is not enough documentation for a curriculum, and the possible speakers are uncontacted, so no programme should be planned in their name. A K-12 unit could treat Katawixi and the Katukinan family as part of Amazonian linguistic diversity and the rights of isolated peoples. The related Kanamari language is better documented.
\end{document}